\section{Phương pháp}
\label{sec:method}

\subsection{Nhánh 1: Mô hình BiLSTM-CRF tự cài đặt từ đầu}
\label{sec:method-bilstm}

Nhánh 1 xây dựng hoàn toàn từ đầu kiến trúc gán nhãn chuỗi kết hợp mạng hồi quy hai chiều và trường ngẫu nhiên có điều kiện (\textbf{BiLSTM-CRF}) trên thư viện PyTorch \citep{lample-etal-2016-neural, ma-hovy-2016-end}. Toàn bộ bộ từ điển (\texttt{vocab.py}), cơ chế đệm động kèm mặt nạ (\texttt{dataset.py}), tầng trích xuất ký tự (\texttt{char\_encoder.py}) và tầng quy hoạch động CRF (\texttt{crf.py}) đều do nhóm tự lập trình nhằm kiểm soát chặt chẽ từng thành phần trong thí nghiệm cắt bỏ (\textit{ablation study}).

\paragraph{Tầng biểu diễn từ và ký tự.}
Với chuỗi đầu vào $\mathbf{x} = (x_1, x_2, \dots, x_n)$, mỗi token $x_i$ được ánh xạ sang vector nhúng từ $\mathbf{e}^{\text{word}}(x_i) \in \mathbb{R}^{d_w}$ ($d_w = 100$) thông qua bảng tra cứu học trực tiếp từ tập huấn luyện, trong đó token đệm \texttt{<PAD>} có chỉ số $0$ (\texttt{padding\_idx=0}) và các từ ngoài từ điển được gán về \texttt{<UNK>}. Để xử lý hiện tượng từ hiếm (OOV) và tận dụng hình thái viết hoa của tên riêng tiếng Việt, vector từ được ghép nối với vector biểu diễn cấp ký tự $\mathbf{e}^{\text{char}}(x_i) \in \mathbb{R}^{d_c}$ trích xuất bởi Char-CNN (tích chập 1D kết hợp \textit{Max-Pooling}) hoặc Char-LSTM:
\begin{equation}
    \mathbf{u}_i = \left[\,\mathbf{e}^{\text{word}}(x_i) \,;\, \mathbf{e}^{\text{char}}(x_i)\,\right] \in \mathbb{R}^{d_w + d_c}.
\end{equation}

\paragraph{Tầng mã hóa ngữ cảnh BiLSTM.}
Chuỗi biểu diễn sau khi đi qua lớp điều chuẩn \textit{Dropout} ($p=0.5$) được nén bằng kỹ thuật \texttt{pack\_padded\_sequence} để triệt tiêu hoàn toàn ảnh hưởng của các vị trí đệm ở cả chiều xuôi và chiều ngược, sau đó đưa qua mạng LSTM hai chiều:
\begin{equation}
    \mathbf{h}_i = \left[\,\overrightarrow{\text{LSTM}}(\mathbf{u}_{1:i}) \,;\, \overleftarrow{\text{LSTM}}(\mathbf{u}_{i:n})\,\right] \in \mathbb{R}^{2d_h},
\end{equation}
với tổng kích thước ẩn $2d_h = 256$. Một tầng chiếu tuyến tính ánh xạ $\mathbf{h}_i$ sang vector điểm phát xạ (\textit{emission scores}) $\mathbf{P}_i \in \mathbb{R}^{|\mathcal{Y}|}$ tương ứng với $|\mathcal{Y}| = 21$ nhãn BIO.

\paragraph{Mô hình cơ sở (Baseline) và Tầng CRF.}
Ở cấu hình cơ sở (\textit{Baseline BiLSTM}), mô hình dự đoán độc lập tại từng bước thời gian và tối ưu hàm mất mát \textit{Masked Cross-Entropy} (bỏ qua vị trí đệm bằng mặt nạ nhị phân \texttt{mask}). Ở cấu hình đầy đủ, tầng \textbf{CRF} bổ sung ma trận chuyển tiếp $\mathbf{A} \in \mathbb{R}^{(|\mathcal{Y}|+2)\times(|\mathcal{Y}|+2)}$ để mô hình hóa ràng buộc hợp lệ giữa các nhãn BIO liên tiếp, tối ưu hàm âm log-hợp lý bằng thuật toán Forward trên miền \textit{log-sum-exp} và giải mã chuỗi tối ưu bằng thuật toán Viterbi.

% --- [Phần của Thành viên 2: Nhánh 2 - PhoBERT & XLM-R sẽ được ghép tại đây] ---
% --- [Phần của Thành viên 3: Nhánh 3 - LLM Few-shot sẽ được ghép tại đây] ---
